%==============================================================
%  lecons/topologie/baire.tex — Théorème de Baire
%==============================================================

\lecontitle{Théorème de Baire}{Topologie — Espaces métriques complets}

%=============================================================
%  § 1 — DÉFINITIONS ET ÉNONCÉ
%=============================================================
\secnum{navyblue}{1}{Espaces complets et théorème de Baire}

\begin{tcolorbox}[thmbox]
  \textbf{Rappels.}\enspace\index{espace métrique complet}
  Un espace métrique $(E,d)$ est \emph{complet} si toute suite de
  Cauchy\index{suite de Cauchy} y converge.
  Une partie $A \subset E$ est \emph{nulle part dense}\index{nulle part dense}
  si son adhérence $\overline{A}$ est d'intérieur vide :
  $\mathring{\overline{A}} = \emptyset$.
  On dit que $A$ est \emph{maigre}\index{ensemble maigre} (ou de première
  catégorie de Baire) si c'est une réunion dénombrable de fermés nulle part
  denses.

  \medskip
  \textbf{Théorème de Baire.}\enspace\index{Baire (théorème de)}
  Soit $(E,d)$ un espace métrique complet non vide. Alors :
  \begin{itemize}
    \item Toute intersection dénombrable d'ouverts denses de $E$ est dense
          dans $E$.
    \item $E$ ne peut pas s'écrire comme réunion dénombrable de fermés
          d'intérieur vide (autrement dit, $E$ n'est pas maigre dans lui-même).
  \end{itemize}
\end{tcolorbox}

\rubriq{Signification intuitive}
Un espace complet est « topologiquement gros » : il est impossible de le
recouvrir par une infinité dénombrable d'ensembles « minces ». Les deux
formulations sont contrapositives l'une de l'autre par passage au complémentaire.
Ce résultat distingue radicalement $\mathbb{R}$ (complet, non maigre) de
$\mathbb{Q}$ (incomplet, maigre dans lui-même).

\decsep{}

%=============================================================
%  § 2 — DÉMONSTRATION
%=============================================================
\secnum{navyblue}{2}{Démonstration}

\rubriq{Idée de la preuve}
On construit récursivement des boules fermées emboîtées, de rayons tendant
vers $0$, en tirant parti de la densité de chaque ouvert. La complétude
garantit que les centres forment une suite de Cauchy qui converge vers un point
appartenant à toutes les boules, et donc à toutes les intersections.

\vspace{10pt}
\noindent{\bfseries\color{navyblue} Preuve.}

\begin{mdframed}[style=proofbar]
  Soient $(U_n)_{n \geq 1}$ des ouverts denses de $E$. Soit $V$ un ouvert
  non vide quelconque de $E$ ; on veut montrer que $V \cap \bigcap_{n \geq 1} U_n
  \neq \emptyset$.

  \medskip
  On construit par récurrence des boules fermées $B_n = \overline{B}(x_n, r_n)$
  satisfaisant :
  \[
    B_n \;\subset\; U_n \cap \mathring{B}_{n-1}, \qquad 0 < r_n \leq \tfrac{1}{n},
  \]
  en posant $B_0 = \overline{B}(x_0, r_0) \subset V$ pour un $x_0 \in V$
  quelconque.

  \medskip
  \textit{Construction.} Supposons $B_n$ construite. Comme $U_{n+1}$ est dense,
  $U_{n+1} \cap \mathring{B}_n$ est un ouvert non vide. On y choisit
  $x_{n+1}$ et $r_{n+1} \leq \min(1/(n+1),\, r_n - d(x_n, x_{n+1}))$
  de sorte que $B_{n+1} \subset U_{n+1} \cap \mathring{B}_n$.

  \medskip
  \textit{Suite de Cauchy.} Pour $p, q \geq n$, les centres $x_p$ et $x_q$
  appartiennent à $B_n$, donc
  \[
    d(x_p, x_q) \;\leq\; 2r_n \;\leq\; \frac{2}{n} \;\xrightarrow{n\to\infty}\; 0.
  \]
  La suite $(x_n)$ est de Cauchy. Par complétude de $E$, elle converge vers
  un $\ell \in E$.

  \medskip
  \textit{Conclusion.} Pour tout $n \geq 1$ et tout $k \geq n$, on a
  $x_k \in B_n$ (fermée), donc $\ell \in B_n \subset U_n$.
  De plus $\ell \in B_0 \subset V$. Ainsi
  \[
    \ell \;\in\; V \;\cap\; \bigcap_{n \geq 1} U_n \;\neq\; \emptyset.
  \]
  $V$ étant arbitraire, $\bigcap_{n \geq 1} U_n$ est dense dans $E$.
  \hfill$\blacksquare$
\end{mdframed}

\rubriq{Équivalence des deux formulations}

\begin{mdframed}[style=proofbar]
  Supposons $E = \bigcup_{n \geq 1} F_n$ avec chaque $F_n$ fermé d'intérieur
  vide. Les $U_n = E \setminus F_n$ sont alors des ouverts denses, et
  \[
    \bigcap_{n \geq 1} U_n = E \setminus \bigcup_{n \geq 1} F_n = \emptyset,
  \]
  ce qui contredit la densité (car $E \neq \emptyset$).
  \hfill$\blacksquare$
\end{mdframed}

\decsep{}

%=============================================================
%  § 3 — APPLICATIONS, CONTRE-EXEMPLES, EXERCICES
%=============================================================
\secnum{exgreen}{3}{Applications, contre-exemples et exercices}

\noindent{\bfseries\color{navyblue} Applications.}

\begin{mdframed}[style=exbar]
  \textbf{1. $\mathbb{R}$ est indénombrable.}
  Si $\mathbb{R} = \{x_n \mid n \geq 1\}$, alors $\mathbb{R}$ serait réunion
  dénombrable des singletons $\{x_n\}$, qui sont des fermés d'intérieur vide.
  Cela contredirait Baire (car $\mathbb{R}$ est complet). Donc $\mathbb{R}$
  est indénombrable.

  \smallskip\noindent
  \textbf{2. Fonctions continues nulle part dérivables.}
  Dans $(\mathcal{C}([0,1]), \|\cdot\|_\infty)$, complet, l'ensemble
  \[
    D_n = \Bigl\{f : \exists\, x,\; \bigl|f(x+h)-f(x)\bigr| \leq n|h|
    \;\text{ pour tout } h\Bigr\}
  \]
  est un fermé d'intérieur vide. Par Baire, $\bigcup D_n \neq \mathcal{C}([0,1])$ :
  il existe des fonctions continues nulle part dérivables.

  \smallskip\noindent
  \textbf{3. Principe de Banach-Steinhaus.}\index{Banach-Steinhaus (théorème de)}
  Soient $E$ de Banach et $(T_n)$ des opérateurs linéaires continus $E\to F$
  tels que $\sup_n \|T_n x\| < +\infty$ pour tout $x$. Baire, appliqué aux
  fermés $F_k = \{x : \sup_n \|T_n x\| \leq k\}$, donne l'existence d'une
  boule sur laquelle ils sont uniformément bornés, d'où $\sup_n \|T_n\| < +\infty$.
\end{mdframed}

\vspace{6pt}
\noindent{\bfseries\color{counterred} Contre-exemples.}

\begin{mdframed}[style=cexbar]
  \textbf{Complétude indispensable.}\enspace%
  $\mathbb{Q}$, muni de la distance usuelle, est non complet et vérifie
  $\mathbb{Q} = \bigcup_{q \in \mathbb{Q}} \{q\}$ : réunion dénombrable de
  singletons fermés d'intérieur vide. Baire est en défaut.

  \smallskip\noindent
  \textbf{Dense $\neq$ grand au sens de la mesure.}\enspace%
  Pour tout $\varepsilon > 0$, l'ensemble
  $G = \bigcup_{n \geq 1} \bigl]q_n - \varepsilon/2^n,\, q_n + \varepsilon/2^n\bigr[$
  (où $(q_n)$ est une énumération de $\mathbb{Q}$) est un ouvert dense
  de mesure de Lebesgue $\leq 2\varepsilon$ : dense au sens de Baire mais
  de mesure arbitrairement petite.
\end{mdframed}

\vspace{6pt}
\noindent{\bfseries\color{goldaccent} Exercices.}

\begin{mdframed}[style=exobar]
  \begin{enumerate}
    \item Montrer que $\mathbb{R}\setminus\mathbb{Q}$ est non maigre dans
          $\mathbb{R}$ (utiliser Baire et le fait que $\mathbb{Q}$ est maigre).

    \item Soit $f : \mathbb{R}\to\mathbb{R}$ telle que pour tout $x$,
          $f$ est continue en $x$ ou dérivable en $x$. Montrer que l'ensemble
          de continuité est un $G_\delta$ dense.

    \item Montrer que tout espace de Banach de dimension infinie ne peut avoir
          de base algébrique dénombrable (appliquer Baire aux sous-espaces
          vectoriels fermés engendrés par les $n$ premiers vecteurs).

    \item Soit $(f_n)$ une suite de fonctions continues $[0,1]\to\mathbb{R}$
          convergeant simplement vers $f$. Montrer que $f$ est continue sur
          un ensemble dense (théorème d'Osgood\index{Osgood (théorème d')}).

    \item Montrer que $\mathbb{R}$ ne peut s'écrire comme réunion dénombrable
          d'ensembles fermés deux à deux disjoints, sauf si l'un d'eux a un
          intérieur non vide.
  \end{enumerate}
\end{mdframed}

\leconfooter{R.~Baire, \textit{Sur les fonctions de variables réelles}, Ann. di Mat.\ (1899)}
